\documentclass[11pt,a4paper]{article}
\usepackage{gtu-el-yazisi}
\begin{document}
\notbilgi{Elektronik Elemanlar}{Yarı iletken ve diyot (s. 1--20)}{Elektronik-elemanlar-notlarim.pdf}

\sayfa{1}
{\small\color{kursunkalem}(Sol üstte okunamayan kısa bir not: \okunamadi{}.)}

\noindent Kitap (ders içeriğine ok ile bağlanmış): \textit{Physics of Semiconductor Devices}, Sze S. M., 2ed edition.

\noindent \textbf{Ders içeriği:} Katı yapılar, Elektron \okunamadi{için} hareketleri, yalıtkan, iletkenler, yarı iletken fiziği.

\medskip
\noindent \textcircled{\small 1} Diyot $\rightarrow$ (PN Junction)\quad
\begin{circuitikz}[scale=0.6, transform shape]\draw (0,0) to[D] (1.6,0);\end{circuitikz}

\noindent \textcircled{\small 2} BJT (Bipolar Junction Transistor)\quad
\begin{circuitikz}[scale=0.45, transform shape]\draw (0,0) node[npn]{};\end{circuitikz}

\noindent \textcircled{\small 3} FET (Field effect transistor) \ {\small\color{kursunkalem}(üstüne küçük harflerle: ``Alan etkili transistör'')}

\noindent \qquad JFET\ {\small\color{kursunkalem}(altında küçük: ``Junction'')}, \qquad MOSFET $\rightarrow$
\begin{circuitikz}[scale=0.45, transform shape]\draw (0,0) node[nmos]{};\end{circuitikz}
\ $\nearrow$ AC, \ $\searrow$ \textcircled{DC}

\section{Kaynaklar}
\noindent Ders kitabı: $\star$ \textit{Microelectronic Circuits} $\rightarrow$ 2016 \quad (Sedra, Smith) \quad
--- The Oxford Series in electrical and computer engineering

\noindent Elektronik $\rightarrow$ Mehmet Sait Türköz $\rightarrow$ Birsen Yayınevi, 2004

\section{Değerlendirme}
\begin{itemize}
  \item 1 ödev ($\sim$\%\okunamadi{5}), 1 proje $\rightarrow$ \textcircled{LTSpice}
  \item Vize $\rightarrow$ \%30
  \item Final $\rightarrow$ \%50
\end{itemize}

\sayfa{2}
{\small\color{kursunkalem}(Sol üst köşede küçük: ``Ders 1'')}
\section{Semiconductors}
\begin{itemize}
  \item Intrinsic Semiconductors (Has yarıiletken)
  \item Doped Semiconductor (katkılı ,,)
\end{itemize}
$\rightarrow$ Elemental Semi $\rightarrow$ \fbox{Silicon}, \ Germanyum\\
$\hookrightarrow$ GaAs $\rightarrow$ Gallium Arsenide, \quad $\rightarrow$ compound semi

\sekilnotu[2]{Silikon kristal kafesi (solda küçük: ``Intrinsic (Has yarıiletken) Silicon''): $3\times3$ düzende içinde $+4$ yazan daireler (iyon çekirdekleri); komşu çekirdekler arasında kovalent bağları gösteren oval çizgiler ve bağların üstünde/altında ``$-$'' (elektron) işaretleri. Sağ üstte (kırmızı) ``Outer shell (dış yörünge)''. Ortadaki çekirdeğin yanında ok ile ``Hole''; sağda ok ile ``covalent \okunamadi{Bonds}''; kırmızı bir çizgi serbest elektrondan ``free electron'' yazısına gidiyor.}
$0^\circ$K'de tüm bağlar \okunamadi{mükemmel} $\rightarrow$ $T\uparrow$

\subsection{Hole (Elektron boşluğu)}
{\small\color{kursunkalem}(Sol kenarda eğik, küçük: ``Boşluk hareketi nasıl oluyor?'' \okunamadi{}.)}

\sekilnotu[2]{Boşluğun hareketi: altı satır halinde küçük daire dizileri (her daire bir elektron, üstünde ``$-$''). Her satırda bir daire eksik (boşluk) ve boşluk satır satır sola kayıyor (1. satırda sağdan ikinci konumda, 6. satırda en solda). İlk satırın üstünde küçük bir ok. Sağda mavi/siyah ok $\rightarrow$ ``elektron hareketi''; altında kırmızı ok $\leftarrow$ ``gerçekteki boşluk hareketi''.}

\subsection{Recombination (Rekombinasyon)}
\sekilnotu[2]{Tek satır daire dizisi, içinde bir boşluk; çizgiyle sağdaki ``free electron'' notuna (üstünde küçük daire içinde elektron işareti) bağlanmış.}
anhilate: serbest elektronun boşlukla yer değiştirmesi

\noindent Termal dengede (Has yarı iletkende)\\
Thermal generation $=$ Recombination\\
\hspace*{1.5em}sıcaklıkla birlikte yeni boş elektron \okunamadi{oluşması}

\sayfa{3}
\sekilnotu[3]{Sayfanın üstünde, çerçevenin dışına taşan enerji bandı çizimi: üstte yatay çizgi $E_c$ (ok ile ``Conduction Band Energy''), altta yatay çizgi $E_v$ (ok ile ``\okunamadi{Valence} Band Energy''); aralarında dikey ok ve ``Bandgap $\rightarrow$ 1,1 eV''.}

\subsection{Free carriers (Serbest Taşıyıcı)}
\begin{itemize}
  \item electrons (Elektron)
  \item holes (Boşluk)
\end{itemize}

\noindent (Has yarı iletken)
\formulkutu{n = p = n_i}
$\rightarrow$ (3,1) $\rightarrow$ sayfa (138)\\
$n$: elektron sayısı, \quad $p$: holes (Boşluk) sayısı

\[ n_i = B\,T^{3/2}\,e^{-E_g/2kT} \]
\begin{itemize}
  \item $E_g$: Band gap Energy (Band aralığı ,,)
  \item $B$: Parametre (silikon) $= 7{,}3\cdot10^{15}\ \mathrm{cm^{-3}K^{-3/2}}$
\end{itemize}
\[ T = 300\ \mathrm{K} \;\Rightarrow\; n_i = 1{,}5\times10^{10}\ \text{carriers/cm}^3 \]
silicon $\rightarrow$ $5\cdot10^{22}$ atom/cm$^3$

\subsection{Doped (katkılı) semiconductor}
\formulkutu{n\cdot p = n_i^2}
$\hookrightarrow$ Has yarıiletken ve katkılı yarıiletken için geçerli.

\sekilnotu[3]{Katkılı silikon kafesi: $3\times3$ düzende $+4$ daireleri ve aralarında kovalent bağ ovalleri; orta satırın ortasındaki çekirdek $+5$ (büyük daire). Üst satırın ortasında bir serbest elektron (daire içinde ``$-$''), ok ile ``Free electron''. $+5$ çekirdeğinden ok ile ``Pentavalent atom'' $\hookrightarrow$ ``impurity atom''. Çizimden aşağıdaki ``n tipi'' kutusuna kavisli bir ok gidiyor.}

\[ p_n \cdot n_n = n_i^2 \]
{\small\color{kursunkalem}(Alt indisler daire içine alınmış: $n$ indisleri ``yarıiletkenin tipi'', $i$ indisi ``intrinsic''.)}

\medskip
\noindent \fbox{n tipi:}

\noindent\fbox{\parbox{0.9\linewidth}{%
$p_n$ : n-tipi yarıiletkende boşluk sayısı\\
$n_n$ : \ ,, \quad ,, \quad elektron sayısı}}

\medskip
\noindent \textbf{Doped semiconductor}
\begin{itemize}
  \item n-tipi yarıiletken: çoğunluk taşıyıcıları elektronlar
  \item p-tipi yarıiletken: çoğunluk taşıyıcı boşluk
\end{itemize}

\sayfa{4}
\subsection{p tipi}
\sekilnotu[4]{Silikon kafesi: $3\times3$ düzende $+4$ daireleri ve kovalent bağ ovalleri; orta satırın ortasındaki çekirdek $+3$. Bu çekirdekten sağa doğru üç ok (etiketlere) ve alttan, boş bağa doğru bir ok çıkıyor.}
\begin{itemize}
  \item Trivalent (Bor)
  \item impurity Atom
  \item $N_A$: Acceptor Atom (+3 valance electron)\\ \hspace*{1em}$N_A$ : p tipi yarı iletken
  \item $N_D$: Donor Atoms (+5 valance electron)\\ \hspace*{1em}$\llcorner$ n tipi yarı iletken
\end{itemize}

\[ n_p \cdot p_p = n_i^2 \]
$p_p$ : p-tipi yarıiletkendeki boşluk sayısı\\
$n_p$ = p-tipi yarıiletkendeki elektron sayısı

\medskip
\noindent\fbox{\parbox{0.9\linewidth}{%
$p_n \cdot n_n = n_i^2$\\[1mm]
\textcircled{$N_D$} $\rightarrow$ \fbox{n-tipi} yarıiletken\\[1mm]
\fbox{$n_n \simeq N_D$}\\[1mm]
$\Big\{\;\fbox{$p_n = \dfrac{n_i^2}{N_D}$}\;\Big\}$ n tipi yarıiletken}}

\medskip
\noindent \fbox{p-tipi yarıiletken}

\noindent \fbox{$p_p \simeq N_A$} $\rightarrow$ Acceptor Atom yoğ.

\[ p_p\, n_p = n_i^2 \;\Rightarrow\; \boxed{n_p \simeq \frac{n_i^2}{N_A}} \quad \text{(Elektron sayısı)} \]

\sayfa{5}
\begin{center}
\begin{tabular}{lcc}
  \toprule
   & n-tipi: & p-tipi: \\ \midrule
  Çoğunluk taşıyıcıları (Majority carriers) & electrons & holes \\
  Azınlık taşıyıcıları (minority carriers) & holes & electrons \\ \bottomrule
\end{tabular}
\end{center}

\begin{ornek}[3.2]
  (sayfa 131) n-type silicon, katkı yoğunluğu (dopant concentration) $N_D = 10^{17}/\mathrm{cm}^3$,
  electron ve boşluk yoğunluğunu $T = 300$ K bulunuz.
\end{ornek}
\begin{cozum}
  \[ n_n = N_D \simeq 10^{17}/\mathrm{cm}^3 \]
  \[ p_n = \frac{n_i^2}{N_D} \;\Rightarrow\; n_i = 1{,}5\cdot10^{10}/\mathrm{cm}^3 \]
  \[ p_n = \frac{(1{,}5\cdot10^{10})^2}{10^{17}} = 2{,}25\cdot10^{3}/\mathrm{cm}^3 \]
\end{cozum}

\section{Current Flow in semiconductors}
(sürükleme akımı)
\subsection{Drift Current}
\[ v_{p\text{-}drift} = \mu_p \cdot E \qquad \text{(Elektrik alanı)} \]
{\small\color{kursunkalem}(Küçük notlar: $v_{p\text{-}drift}$ için ``hızı''; $\mu_p$ için ``mobilite hole''.)}
\[ \mu_p = 480\ \mathrm{cm^2/V\!\cdot\!s} \]
\[ v = \mu\cdot E \quad \text{(hızı; Elektrik alanı)}, \qquad
   \frac{\mathrm{cm}}{\mathrm{s}} = \mu\cdot\frac{\mathrm{V}}{\mathrm{cm}} \;\Rightarrow\; \frac{\mathrm{cm}^2}{\mathrm{V\cdot s}} \]
\[ v_{n\,drift} = -\mu_n \cdot E \]
Elektron mobilitesi $\hookrightarrow$ $1350\ \mathrm{cm^2/V\!\cdot\!s}$

\sayfa{6}
\sekilnotu[6]{Üç boyutlu dikdörtgen prizma (çubuk); üstünde sağa doğru ok ``$\rightarrow E$''. Ön yüzde ``$+$'' $\rightarrow$ ``Hole'' (sağa ok) ve ``$-$'' $\leftarrow$ ``electrons'' (sola ok).}

\noindent in one second:
\[ I_p = A\,q\,p\,v_{p\,drift} \]
\[ A\cdot q\cdot p\cdot v \quad (A:\ \mathrm{cm}^2,\ q:\ \mathrm{C},\ p:\ 1/\mathrm{cm}^3,\ v:\ \mathrm{cm/s}) \;\rightarrow\; \mathrm{C} \]

\[ I_p = A\,q\,p\,\mu_p\,E \]
{\small\color{kursunkalem}(Soldaki küçük not: ``\okunamadi{deliklerin}''.)}

$\rightarrow$ $J$ $\rightarrow$ current density
\[ J_p = \frac{I_p}{A} = q\cdot p\cdot\mu_p\cdot E \]
{\small\color{kursunkalem}(Soldaki küçük not: ``elektronların''.)}
\[ I_n = -A\cdot q\cdot n\cdot v_{n\,drift} \qquad (v_{n\,drift} = -\mu_n E) \]
\[ J_n = \frac{I_n}{A} = q\cdot n\cdot\mu_n\cdot E \]
\[ J = J_p + J_n = q\,(p\mu_p + n\mu_n)\,E \]
\[ J = \sigma\cdot E \qquad \hookrightarrow \text{conductivity (iletkenlik)} \]
\[ \rho = \frac{1}{\sigma} = \frac{1}{q\,(p\mu_p + n\mu_n)} \qquad (\rho:\ \text{resistivity}) \]
\[ \rho = \frac{E}{J}, \qquad \rho = R\,\frac{A}{\ell} \quad (A:\ \text{kesit alanı},\ \ell:\ \text{uzunluk}) \]

\sayfa{7}
\begin{ornek}[3.3]
  Resistivity $=$ ?\\
  a) intrinsic silicon\\
  b) p-type \ ,, $\rightarrow$ $N_A = 10^{16}/\mathrm{cm}^3$\\[1mm]
  $n_i = 1{,}5\cdot10^{10}/\mathrm{cm}^3$\\
  has yarıiletken $\rightarrow$ $\mu_n = 1350\ \mathrm{cm^2/V\!\cdot\!s}$, \ $\mu_p = 470\ \mathrm{cm^2/V\!\cdot\!s}$\\
  katkılı yarıiletken $\rightarrow$ $\mu_n = 1110\ \mathrm{cm^2/V\!\cdot\!s}$, \ $\mu_p = 400\ \mathrm{cm^2/V\!\cdot\!s}$
\end{ornek}
\begin{cozum}
  \textbf{a)}
  \[ p = n = n_i = 1{,}5\cdot10^{10}/\mathrm{cm}^3 \]
  \[ \rho = \frac{1}{q\,(p\mu_p + n\mu_n)} = \frac{1}{1{,}6\cdot10^{-19}\,(1{,}5\cdot10^{10}\cdot 470 + 1{,}5\cdot10^{10}\cdot1350)} \]
  \[ = 2{,}28\times10^{\okunamadi{5}}\ \Omega\,\mathrm{cm} \;\rightarrow\; \text{Has yarıiletken} \]
  \textbf{b)} p tipi yarıiletken
  \[ p_p \simeq N_A = 10^{16} \]
  \[ n_p = \frac{n_i^2}{N_A} = \frac{(1{,}5\cdot10^{10})^2}{10^{16}} = 2{,}25\cdot10^{4}/\mathrm{cm}^3 \]
  \[ \rho = \frac{1}{q\,(p\mu_p + n\mu_n)} = \frac{1}{1{,}6\times10^{-19}\,(10^{16}\times400 + 2{,}25\times10^{4}\times1110)} \]
  \noindent $\displaystyle \simeq \frac{1}{1{,}6\times10^{-19}\times10^{16}\times400} = 1{,}66\ \Omega\,\mathrm{cm}$\footnote{Claude notu: $1/0{,}64 \approx 1{,}56$ olduğundan muhtemelen $1{,}56\ \Omega\,\mathrm{cm}$ olmalı.}
\end{cozum}

\sayfa{8}
\section{Diffusion Current}
\sekilnotu[8]{Üç boyutlu çubuk; sol uca ok ile ``Hole injection'', çubuğun içinde çok sayıda ``$+$'' işareti; altında sağa doğru ok ve ``$x$''.}
\begin{center}
\begin{tikzpicture}[line width=0.7pt]
  \draw[->] (-0.3,0) -- (3.4,0) node[right] {$x$};
  \draw[->] (0,0) -- (0,2.2) node[above] {$p(x)$};
  \draw[murekkep,domain=0:3.2,smooth,variable=\x] plot ({\x},{1.9*exp(-1.1*\x)});
  \node[left,align=right,font=\small] at (0,1.2) {Hole\\ concentration};
\end{tikzpicture}
\end{center}

\formulkutu{J_p = -q\,D_p\,\frac{dp(x)}{dx}}
$D_p$: difüzyon sabiti (Boşluk)
\[ J_n = q\,D_n\,\frac{dn(x)}{dx} \qquad (D_n:\ \text{elektron difüzyon sabiti}) \]

\begin{center}
\begin{tikzpicture}[line width=0.7pt]
  \draw[->] (0,0) -- (3.4,0);
  \draw[->] (0,0) -- (0,2.2) node[above] {$n(x)$};
  \draw[murekkep,domain=0:3.2,smooth,variable=\x] plot ({\x},{1.9*exp(-1.3*\x)});
  \node[left,align=right,font=\small] at (0,1.2) {Electron\\ concentration};
\end{tikzpicture}
\end{center}
Diffusion constants $\hookrightarrow$ $D_p = 12\ \mathrm{cm^2/s}$, \quad $D_n = 35\ \mathrm{cm^2/s}$

\begin{ornek}[3.4]
  $p(x) = p_0\,e^{-x/L_p}$\\
  Find the hole density at $x = 0$.\\
  $p_0 = 10^{16}/\mathrm{cm}^3$, \ $L_p = 1\,\mu$m, \ $D_p = 12\ \mathrm{cm^2/s}$\\
  $I_p = {}$? \quad (Bar cross sectional Area $10\ \mu\mathrm{m}^2$)
\end{ornek}
\begin{cozum}
  \begin{align*}
    J_p &= -q\,D_p\,\frac{dp(x)}{dx} \\
        &= -q\,D_p\,\frac{d}{dx}\left[p_0\,e^{-x/L_p}\right] \\
        &= q\,\frac{D_p}{L_p}\,p_0\,e^{-x/L_p} \\
    J_p(0) &= q\,\frac{D_p}{L_p}\,p_0 \\
        &= 1{,}6\cdot10^{-19}\cdot\frac{12}{1\cdot10^{-4}}\cdot10^{16} = 192\ \mathrm{A/cm^2}
  \end{align*}
  \[ I_p = J_p\times A \;\Rightarrow\; 192\cdot10^{-7} = 19{,}2\ \mu\mathrm{A} \]
\end{cozum}

\sayfa{9}
\section{Relationship between $D$ and $\mu$}
\noindent \fbox{Einstein relationship}
\[ \frac{D_n}{\mu_n} = \frac{D_p}{\mu_p} = V_T \]
$V_T$ $\rightarrow$ Termal Gerilim, Thermal Voltage $\rightarrow$
\[ V_T = \frac{kT}{q} \]
{\small\color{kursunkalem}(Çok küçük, silik notlar ok ile: $k$ için ``Boltzmann sabiti'', $T$ için ``sıcaklık'', $q$ için ``\okunamadi{yük}''.)}
\[ T = 300\ \mathrm{K} \;\Rightarrow\; V_T \simeq 25{,}9\ \mathrm{mV} \]
\[ \simeq 26\ \mathrm{mV},\ 25\ \mathrm{mV} \]

\sayfa{10}
{\small\color{kursunkalem}(Sol üst köşede küçük: ``2. Ders'')}
\section{PN Junction}
$n_i = 1{,}5\cdot10^{10}$ silikonda \qquad\qquad $n\cdot p = n_i^2$

\sekilnotu[10]{Üç gösterim. (1) Dikdörtgenin sol yarısı ``P-type silicon'', sağ yarısı ``n-type silicon''; her iki uçta kalın metal kontak; sol uçtan ok ile ``Anode'' ve ``metal kontakt'' notu, sağ uçtan ok ile ``Cathode''. (2) Daha küçük kutu: ``Anode --- P | N --- Cathode''. (3) Altında, iki yanında uzantı çizgileri olan diyot sembolü.}

\subsection{Open circuit}
\sekilnotu[10]{Açık devre PN eklemi: büyük dikdörtgen; sol bölge (köşesinde ``P'') içinde dağınık ``$+$'' işaretleri; ortada dikey iki sütun halinde daire içine alınmış işaretler (solda ``$-$'', sağda ``$+$''; tükenme bölgesi, üstünde süslü parantez); sağ bölge (n) içinde dağınık ``$-$'' işaretleri ve birkaç ``$+$''. Üstte oklar: ``holes diff'' $\rightarrow$ (iki kez), $\leftarrow$ ``electrons diff'' ve altı çizili ``$I_{Diff}$''. Sağ kenarda oval içinde ``uncovered Atom'' ve ona giden uzun, kavisli bir ok.}
\begin{itemize}
  \item Depletion Region (Fakirleşmiş bölge) (Burada taşıyıcı yok) (recombination olur)
  \item Sol bölgenin altında: Majority carriers (çoğunluk T); ilk önce difüzyon gerçekleşir
  \item Ortada, ok ile $E$: (Elektrik alan oluşur (difüzyon))
  \item Sağ bölgenin altında: Çoğunluk T
  \item Sağda: $\leftarrow$ $I_S$ (drift current); $\leftarrow$ holes drift; electrons drift $\rightarrow$
  \item Band Charges (Hareketsizdir, sabittir) \quad (ok ile tükenme bölgesindeki yüklere bağlı)
\end{itemize}

\noindent Artıyükler elektrik alanla aynı yönde, eksiler ise ters yönde.

\noindent Dengede $\Rightarrow$ \quad $I_D = I_S$ \quad ($I_D$: difüzyon akımı; \ $I_S$: drift current)

\noindent {\small\color{kursunkalem}Elektrik alan neticesinde}
\begin{center}
\begin{tikzpicture}[line width=0.7pt]
  \draw[->] (0,0) -- (4.6,0) node[right] {$x$};
  \draw[->] (0,0) -- (0,2.5);
  \node[rotate=90] at (-0.35,1.25) {Potential};
  \draw[murekkep] (0,0.1) -- (1.6,0.1) .. controls (2.2,0.1) and (2.1,1.9) .. (2.8,1.9) -- (3.6,1.9);
  \draw[<->] (3.6,0.1) -- (3.6,1.9);
  \node[right,align=left] at (3.65,1.0) {Barrier\\ Voltage\\ ($V_0$)};
\end{tikzpicture}
\end{center}

\[ V_0 = V_T\,\ln\!\left(\frac{N_A\,N_D}{n_i^2}\right) \]
$N_A$: Acceptor atom (+3) p tipi; \quad $N_D$: Donor atom (+5) n tipi

\sayfa{11}
\section{Farklı katkı yoğunluklarında pn jonksiyonu}
\noindent $N_A > N_D$

\sekilnotu[11]{Dikdörtgen PN eklemi: sol (p) bölgede çoğunlukla ``$+$'' ve birkaç ``$-$'' işareti; ortada dar bir sütun halinde 6 daire içinde ``$-$'' ($-x_p$ ile $0$ arası), onun yanında daha geniş bir bölgede iki sütun halinde daire içinde ``$+$'' ($0$ ile $x_n$ arası); sağ (n) bölgede çoğunlukla ``$-$'' ve birkaç ``$+$''. Altında $-x_p$, $0$, $x_n$ işaretleri; sol altta küçük ``depletion \okunamadi{region}'' notu; $N_A>N_D$ yazısının üstünde çok küçük, okunamayan notlar. Sağda çerçeve içinde $x_n/x_p = N_A/N_D$ (ok ile $x_n$'e küçük bir not: \okunamadi{uzunluk}).}

\[ \left.\begin{aligned} |Q_+| &= q\,A\,x_n\,N_D \\ |Q_-| &= q\,A\,x_p\,N_A \end{aligned}\right\}\;\rightarrow\; q\,A\,x_n\,N_D = q\,A\,x_p\,N_A \]
\formulkutu{\frac{x_n}{x_p} = \frac{N_A}{N_D}}

\noindent $x_n$ : n tipi bölgedeki fakirleşmiş bölgenin uzunluğu\\
$x_p$ : p tipi \ ,, \ ,, \ ,, \ ,,

\[ W = x_n + x_p = \sqrt{\frac{2\varepsilon_s}{q}\left(\frac{1}{N_A}+\frac{1}{N_D}\right)V_0} \qquad (W:\ \text{the width of depletion region}) \]
\[ \varepsilon_s = \varepsilon_r\cdot\varepsilon_0 = (11{,}7)\,(8{,}85\cdot10^{-14}\ \mathrm{F/cm}) \]

\noindent $\displaystyle \frac{x_n}{x_p} = \frac{N_A}{N_D} \;\Rightarrow\; x_p = x_n\,\frac{N_A}{N_D}$\footnote{Claude notu: muhtemelen $x_p = x_n\,N_D/N_A$ olmalı; sağdaki türetmede de böyle kullanılmış.}
\hfill (sağdaki türetme:)
\begin{align*}
  x_n + x_p &= W \\
  x_n + x_n\,\frac{N_D}{N_A} &= W \\
  x_n\left(1 + \frac{N_D}{N_A}\right) &= W \\
  \boxed{x_n = \frac{N_A}{N_A + N_D}\,W} & \\
  \boxed{x_p = W\,\frac{N_D}{N_A + N_D}} &
\end{align*}

\[ |Q_+| = |Q_-| = Q_J \]
\[ Q_J = A\,q\left(\frac{N_A N_D}{N_A + N_D}\right)W \]
\formulkutu{Q_J = A\sqrt{2\varepsilon_s\,q\left(\frac{N_A N_D}{N_A + N_D}\right)V_0}}

\sayfa{12}
\section{The PN junction with an Applied Voltage}
\sekilnotu[12]{Sol üst: ``Open circuit'' (kutulu başlık). P|N dikdörtgeni, uçlarda kalın metal kontaklar; ortada daire içinde ``$-$'' (sol sütun) ve ``$+$'' (sağ sütun) işaretlerinden oluşan dar tükenme bölgesi. Altında oklar $\rightarrow I_D$ ve $\leftarrow I_S$. Daha altta potansiyel basamağı (S biçimli eğri), yüksekliği $V_0$.}
\sekilnotu[12]{Sağ üst: ``Reverse bias'' (kutulu başlık). Aynı yapı; üstünde dış kaynak $V_R$ (iki uca dolanan teller), tükenme bölgesi daha geniş (solda $-$, sağda $+$ daireleri birbirinden uzak). Altında $\rightarrow I_D$ ve $\leftarrow I_S$. Potansiyel basamağı, yüksekliği $(V_0 + V_R)$.}
\sekilnotu[12]{Sol alt: ``Forward bias'' (kutulu başlık). Aynı yapı; üstte dış kaynak $V_F$, tükenme bölgesi daha dar (daire sütunları birbirine yakın). Altında uzun ok $\rightarrow I_D$ ve kısa $\leftarrow I_S$. Potansiyel basamağı, yüksekliği $(V_0 - V_F)$.}

\subsection{Reverse}
\[ W = x_n + x_p = \sqrt{\frac{2\varepsilon_s}{q}\left(\frac{1}{N_A}+\frac{1}{N_D}\right)(V_0 + V_R)} \]
\[ Q_J = A\sqrt{2\varepsilon_s\,q\left(\frac{N_A N_D}{N_A+N_D}\right)(V_0 + V_R)} \]

\noindent Reverse $\rightarrow$ $I = I_S$

\noindent Forward $\rightarrow$ \fbox{$I = I_S\left(e^{V/V_T} - 1\right)$}

\bigskip
\sekilnotu[12]{Sağ altta diyot sembolü (anot solda), üstünde sola doğru küçük ok; solda ``$+$'' (yanında kırmızı, okunamayan küçük not), sağda ``$-$''; altında ``$I_S$'' yazısı.}
\noindent \fbox{$I_D = I_S\,e^{V_D/nV_T}$} \qquad $V_T = \dfrac{kT}{q}$ \ \okunamadi{$\approx 25$ mV}

\medskip
\noindent $I_S$ sıcaklıktan çok etkilenir

\sayfa{13}
\section{The PN Junction with applied Voltage}
{\small\color{kursunkalem}(Sol kenarda küçük bir yıldız işareti.)}
\begin{center}
\begin{tikzpicture}[line width=0.7pt, scale=0.95]
  \draw[->] (-3.2,0) -- (3.4,0) node[right] {$V$};
  \draw[->] (0,-1.7) -- (0,2.7) node[above] {$I$};
  \draw[murekkep] (-2.6,-1.6) .. controls (-2.5,-0.5) and (-2.45,-0.05) .. (-2.1,0) -- (0.4,0)
        .. controls (0.9,0.05) and (1.1,0.6) .. (1.3,1.4) .. controls (1.4,1.8) and (1.5,2.1) .. (1.75,2.5);
  \node[below,font=\scriptsize] at (-2.4,0) {$-10$ V};
  \node[below,font=\scriptsize] at (-1.2,0) {$-5$ V};
  \node[below,font=\scriptsize] at (0.65,0) {$0{,}7$};
  \node[right,font=\small,align=left] at (2.0,2.3) {Standart Diyot (silikon)};
\end{tikzpicture}
\end{center}

\sekilnotu[13]{Grafiğin altında üç diyot gösterimi: (1) diyot sembolü, solda ``$V_1$'' ($+$), sağda ``$V_2$'' ($-$), üstünde küçük bir ok; (2) iki ucunda nokta olan ``P | n'' kutusu; (3) silindirik fiziksel diyot (sağ ucuna yakın kalın bant), sol ucunda ``$+$'', sağ ucunda ``$-$''; üstünde ``ileri'' ve ``$V_1 > V_2$''.}

\noindent \fbox{Reverse Breakdown}
\begin{itemize}
  \item[$\rightarrow$] Zenner effect $\rightarrow$ elektrik alanın etkisiyle bağların kopması
  \item[$\rightarrow$] Avalanche effect
\end{itemize}

\sayfa{14}
\section{Capacitive effects}
\begin{itemize}
  \item[$\rightarrow$] \fbox{Depletion} or junction capacitance $\rightarrow$ Reverse region
\end{itemize}
\[ Q_J = A\sqrt{2\varepsilon_s\,q\,\frac{N_A N_D}{N_A+N_D}\,(V_0 + V_R)} \]
\[ Q_J = \alpha\sqrt{V_0 + V_R} \]
\[ \alpha = A\sqrt{2\varepsilon_s\,q\,\frac{N_A N_D}{N_A+N_D}} \]

\begin{center}
\begin{tikzpicture}[line width=0.7pt]
  \draw[->] (0,0) -- (4.2,0) node[below] {$V_R$};
  \draw[->] (0,0) -- (0,2.6) node[left] {$Q_J$};
  \draw[murekkep] (0,0.2) .. controls (0.1,1.4) and (0.6,1.9) .. (3.0,2.2);
  \draw[kursunkalem,dashed] (1.3,0) -- (1.3,2.05);
  \node[below,font=\small] at (1.3,0) {$V_R$};
  \node[font=\small] at (1.9,2.65) {slope $= C_j$};
\end{tikzpicture}
\end{center}

\[ C_j = \left.\frac{dQ_J}{dV_R}\right|_{V_R = V_0} \]
\[ \boxed{C_j = \frac{\alpha}{2\sqrt{V_0 + V_R}}} \;\longrightarrow\; C_j = \frac{C_{j0}}{\sqrt{1 + \dfrac{V_R}{V_0}}} \qquad (C_{j0}:\ V_R = 0) \]
{\small\color{kursunkalem}(Sağ kenarda küçük: ``sayfa 169'')}

\medskip
\noindent \textbf{Difüzyon kapasitesi (Diffusion capacitance)} \hfill $\rightarrow$ Forward Bias, minority carriers in effect
\[ C_d = \left(\frac{\tau_T}{V_T}\right) I \qquad (C_d:\ \text{incremental diffusion capacitance};\ \tau_T:\ \text{mean transit time}) \]
$\rightarrow$ sayfa 169--170 \ Şekil \okunamadi{ok.}

\sayfa{15}
{\small\color{kursunkalem}(Sol üst köşede küçük: ``Ders 3'')}
\section{The ideal Diode}
\sekilnotu[15]{Sol üst: ideal diyot sembolü, üstünde ``(ideal)'' ve akım oku; solda ``Anode'', sağda ``Cathode'', altında $+\ V\ -$. Altında $i$--$V$ eksenleri: $V$ ekseninin solu ``Reverse bias'' (sola ok), sağı ``Forward Bias'' (sağa ok); karakteristik eğrisi çizilmemiş.}
\sekilnotu[15]{Sağ üst: iki anahtar modeli. (1) Açık anahtar, $+\ \ -$ uçları ve $i$ oku: $v<0 \Rightarrow i=0$, ``Reverse (off) (Geri yön) (kesim)''. (2) Kapalı anahtar (kısa devre): $i>0 \Rightarrow v=0$, ``Forward (on) (ileri yön) (iletim)''.}

\noindent
\begin{minipage}[t]{0.47\linewidth}
\centering
\begin{circuitikz}[scale=0.8, transform shape]
  \draw (0,3) node[above] {10V} to[R=$1\,\mathrm{k}\Omega$] (0,1.5) to[D*, l_=$0$ V, i_=$10$ mA] (0,0) node[ground]{};
\end{circuitikz}\\
iletim (on)
\end{minipage}\hfill
\begin{minipage}[t]{0.47\linewidth}
\centering
\begin{circuitikz}[scale=0.8, transform shape]
  \draw (0,3) node[above] {10V} to[R=$1\,\mathrm{k}\Omega$] (0,1.5) to[D, invert, i_=$0$ A] (0,0) node[ground]{};
\end{circuitikz}\\
kesim (off)
\end{minipage}

\bigskip
\noindent \textbf{Yarım dalga doğrultucu (ideal diyot)}
\begin{center}
\begin{circuitikz}[scale=0.9, transform shape]
  \draw (0,0) to[V] (0,2.2) to[D, l=(ideal)] (3,2.2) to[short, -o] (4.4,2.2) node[right] {$+$};
  \draw (3,2.2) to[R=$R$] (3,0) to[short, -o] (4.4,0) node[right] {$-$};
  \draw (0,0) -- (3,0);
  \draw (1.5,0) node[ground] {};
\end{circuitikz}
\end{center}

\begin{center}
\begin{tikzpicture}[line width=0.7pt]
  \draw[->] (-0.2,0) -- (6.4,0) node[right] {$t$};
  \draw[->] (0,-1.2) -- (0,1.5) node[above] {$v_S$};
  \draw[murekkep,domain=0:6.0,smooth,variable=\x,samples=60] plot ({\x},{1.1*sin(120*\x)});
  \draw[kursunkalem,dashed] (0,1.1) -- (5.6,1.1);
  \node[left] at (0,1.1) {$V_p$};
\end{tikzpicture}
\end{center}

\noindent \underline{1. durum:} Giriş kaynağı pozitifken
\sekilnotu[15]{Eşdeğer devre: kaynak ($+$ üstte), üst koldaki diyot yerine iki küçük halka arasında düz tel (kısa devre), yük direnci; çıkış ucuna ``$v_O = v_I$'', tele ``$V_D = 0$ (kısa devre) short circuit''.}

\noindent \underline{Giriş kaynağı negatifken}
\sekilnotu[15]{Eşdeğer devre: kaynak ($-$ üstte), üst kolda ``$+\ v_D\ -$'' ile gösterilen açık devre (iki halka arası boş), akım $i_D = 0$, yük direnci ve $+/-$ çıkış uçları.}

\begin{center}
\begin{tikzpicture}[line width=0.7pt]
  \draw[->] (-0.2,0) -- (6.4,0) node[right] {$t$};
  \draw[->] (0,-0.5) -- (0,1.5) node[above] {$v_O$};
  \foreach \a in {0, 3}{
    \draw[murekkep,domain=0:1.5,smooth,variable=\x,samples=30] plot ({\a+\x},{1.1*sin(120*\x)});
  }
\end{tikzpicture}
\end{center}

\sayfa{16}
\begin{ornek}[4.1]
  \begin{center}
  \begin{circuitikz}[scale=0.85, transform shape]
    \draw (0,0) to[V, l_=$V_S$] (0,2.4) to[R=$100\,\Omega$] (2.2,2.4) to[D, l=(ideal)] (4,2.4)
          to[battery2, l=$12$ V] (4,0) -- (0,0);
  \end{circuitikz}
  \end{center}
  $V_S$: tepe değeri $24$ V olan sinüs. \quad a) Peak value of diode current\\
  b) Maximum reverse-bias voltage of the diode?\\
  c) Fraction of each cycle during which the diode conducts (periyodun kaçında diyot iletimdedir?)
\end{ornek}
\begin{cozum}
  \textbf{c)}
  \begin{center}
  \begin{tikzpicture}[line width=0.7pt]
    \draw[->] (-0.3,0) -- (5.2,0) node[right] {$t$};
    \draw[->] (0,-0.4) -- (0,2.6) node[above] {$V_S$ (24 V)};
    \draw[murekkep,domain=0:3.0,smooth,variable=\x,samples=40] plot ({\x},{2.2*sin(60*\x)});
    \draw[kursunkalem,dashed] (0,1.1) -- (4.4,1.1);
    \node[left] at (0,1.1) {12 V};
    \draw[kursunkalem,dashed] (0.5,0) -- (0.5,1.1);
    \draw[kursunkalem,dashed] (1.5,0) -- (1.5,2.2);
    \draw[kursunkalem,dashed] (2.5,0) -- (2.5,1.1);
    \node[below] at (0.5,0) {$30^\circ$};
    \node[below] at (1.5,0) {$\pi/2$};
    \node[below] at (2.5,0) {$150^\circ$};
    \draw (0.5,-0.7) .. controls (1.1,-1.0) and (1.9,-1.0) .. (2.5,-0.7);
  \end{tikzpicture}
  \end{center}
  \[ 150^\circ - 30^\circ = 120^\circ \;\rightarrow\; 1/3\ \text{cycle} \]
  \textbf{a)}
  \[ I_D = \frac{24 - 12}{100} = 0{,}12\ \mathrm{A} \]
  \textbf{b)} Diyot sembolü (anot solda): sol uçta $-24$, sağ uçta $12$ V.
  \[ -24 - 12 = -36\ \mathrm{V} \qquad \text{(uçlar arasına düşen negatif gerilim)} \]
\end{cozum}

\noindent \fbox{Lojik kapılar}

\sayfa{17}
\begin{ornek}[4.2]
  \textbf{a)} \sekilnotu[17]{Devre (a): ``$D_1$'' diyodunun anodu toprağa bağlı, katodu aşağıda $B$ düğümüne iniyor. Sağda $+10$ V $\rightarrow$ $10\,\mathrm{k}\Omega$ $\rightarrow$ çıkış düğümü $V$ (altında küçük toprak işareti) $\rightarrow$ ``$D_2$'' diyodu (anot üstte), katodu $B$ düğümüne bağlı. $B$ düğümünden $5\,\mathrm{k}\Omega$ ile $-10$ V'a.}
\end{ornek}
\begin{cozum}
  {\small\color{kursunkalem}(Sol kenarda: ``Çözmek için varsayımda bulunalım'')}\\
  $D_1$ is on ve $D_2$ is on \ ($\checkmark$ işaretli)
  \sekilnotu[17]{Çözüm devresi: $+10$ V, $10\,\mathrm{k}\Omega$ üzerinde $1$ mA aşağı; $V$ düğümü $(0\,\mathrm{V})$; $D_2$ üzerinde $1$ mA; $D_1$ üzerinde $1$ mA, $D_1$'in anodu $(0\,\mathrm{V})$ toprakta; $B$ düğümü $(0\,\mathrm{V})$, $D_2$ katodundan $B$'ye ok; $5\,\mathrm{k}\Omega$ üzerinde $2$ mA aşağı, ucu $-10$ V.}
  \noindent iletimde akım akar ve \okunamadi{diyot} kısa devre gibi davranır.
\end{cozum}

\begin{ornek}[4.2\,(b)]
  \sekilnotu[17]{Devre (b): $+10$ V $\rightarrow$ $5\,\mathrm{k}\Omega$ $\rightarrow$ çıkış düğümü $V$ (altında $(0)$; sağında $+V$ ve toprağa $-$ ucu) $\rightarrow$ $D_2$ (akım oku aşağı); $D_1$ anodu toprakta, katodu $B$ düğümünde; $D_2$ katodu $B$'ye bağlı; $B$ düğümü $(0\,\mathrm{V})$, ondan $10\,\mathrm{k}\Omega$ ile $-10$ V'a, üzerinde $1$ mA aşağı. $D_1$'in yanında ``sağlanmaz'' ve daire içinde \okunamadi{$-1$ mA}.}
\end{ornek}
\begin{cozum}
  Varsayım: $D_1$ is on and $D_2$ is on $\Rightarrow$ varsayım yanlış\\
  $D_1$ ve $D_2$ kesimde:
  \sekilnotu[17]{Eşdeğer devre: toprağa bağlı açık uç (boşta), $+10$ V $\rightarrow$ $5\,\mathrm{k}\Omega$ $\rightarrow$ $V$ düğümü, aşağıda iki açık uç $-10$ V ve $-10$ V olarak işaretli, ortadan $10\,\mathrm{k}\Omega$ ile $-10$ V'a. Yanında ``$\Delta V > 0$'', ``hipotez yanlış''.}
\end{cozum}

\sayfa{18}
\noindent $D_1$ is off and $D_2$ is on \ ($\checkmark$ işaretli)
\sekilnotu[18]{Sol devre: $D_1$ anodu toprakta, yanında ``off''; $D_2$ yanında ``on'' ve $1{,}33$ mA aşağı ok. Sağda $+10$ V $\rightarrow$ $5\,\mathrm{k}\Omega$ $\rightarrow$ çıkış düğümü $V$ $\rightarrow$ $D_2$; $D_1$ ve $D_2$ katodları $B$ düğümünde, $B$ düğümü daire içinde ``$3{,}3$ V'' (muhtemelen $3{,}33$ V); $B$'den $10\,\mathrm{k}\Omega$ ile $-10$ V'a, üzerinde $1{,}33$ mA aşağı. Altında $\Delta V = \okunamadi{-3{,}33}$.}
$\rightarrow$
\sekilnotu[18]{Eşdeğer devre: $+10$ V $\rightarrow$ $5\,\mathrm{k}\Omega$ $\rightarrow$ iki halka arasında kapalı anahtar ($D_2$ kısa devre) $\rightarrow$ $B$ düğümü $\rightarrow$ $10\,\mathrm{k}\Omega$ $\rightarrow$ $-10$ V.}
\[ I = \frac{20\,\mathrm{V}}{15\,\mathrm{k}} = 1{,}33\ \mathrm{mA} \]

\bigskip
\noindent \textbf{\textcolor{kirmizikalem}{Exercise 4.6}} \ (sf. 183)
\sekilnotu[18]{Dört devre (a)--(d), hepsi ideal diyot + direnç + çıkış düğümü $V$; her birinin yanında daire içine alınmış (soluk) bir durum notu. (a) $+5$ V $\rightarrow$ $2{,}5\,\mathrm{k}\Omega$ $\rightarrow$ $V$ $\rightarrow$ aşağı bakan diyot $\rightarrow$ toprak; not: ``on''. (b) $+5$ V $\rightarrow$ $2{,}5\,\mathrm{k}\Omega$ $\rightarrow$ $V$ $\rightarrow$ yukarı bakan diyot (anot topraktan yana) $\rightarrow$ toprak; not: ``off''. (c) üstte toprak $\rightarrow$ yukarı bakan diyot $\rightarrow$ $V$ düğümü $\rightarrow$ $2{,}5\,\mathrm{k}\Omega$ $\rightarrow$ $-5$ V; not: ``off''. (d) üstte toprak $\rightarrow$ aşağı bakan diyot $\rightarrow$ $V$ düğümü $\rightarrow$ $2{,}5\,\mathrm{k}\Omega$ $\rightarrow$ $-5$ V; not: ``on''.}

\noindent {\small\color{kursunkalem}(Fotoğrafın sağ kenarında bir sonraki sayfanın yarısı da görünüyor; içeriği s.~19'da aktarılmıştır.)}

\sayfa{19}
\section{4.2 Terminal Characteristics of Junction Diodes}
\noindent
\begin{minipage}[c]{0.56\linewidth}
\centering
\begin{tikzpicture}[line width=0.7pt, scale=0.9]
  \draw[<->] (-3.0,0) -- (3.3,0);
  \draw[<->] (0,-2.0) -- (0,2.5);
  \draw[murekkep] (-2.4,-2.0) .. controls (-2.35,-0.6) and (-2.3,-0.05) .. (-1.9,0) -- (0.9,0)
        .. controls (1.15,0) and (1.3,0.3) .. (1.4,1.0) -- (1.5,2.3);
  \node[below,font=\scriptsize] at (0.95,0) {0,5 V};
  \node[below right,font=\scriptsize] at (1.4,0) {0,7 V};
  \node[left,font=\scriptsize] at (0,0.35) {$I_S$};
  \node[above,font=\scriptsize] at (-2.0,0.05) {$-V_{ZK}$};
  \draw[->] (0,1.9) -- (3.0,1.9) node[above,font=\scriptsize,pos=0.55] {Forward};
  \draw[<->] (-2.9,-0.55) -- (-2.2,-0.55);
  \node[below,font=\scriptsize,align=center] at (-2.6,-0.55) {break\\ down};
  \draw[->] (-2.2,-0.55) -- (-0.1,-0.55) node[above,font=\scriptsize,pos=0.5] {reverse};
\end{tikzpicture}
\end{minipage}\hfill
\begin{minipage}[c]{0.40\linewidth}
\textbf{The Forward-Bias Region}
\[ i = I_S\left(e^{v/V_T} - 1\right) \]
\centering
\begin{circuitikz}[scale=0.7, transform shape]
  \draw (0,0) to[D, o-o, v=$\,$, i=$i$] (3,0);
  \node[above,font=\small] at (0,0.05) {$+$};
  \node[above,font=\small] at (3,0.05) {$-$};
\end{circuitikz}
\end{minipage}

\[ i \simeq I_S\,e^{v/V_T} \qquad\qquad V = V_T\,\ln\frac{i}{I_S} \]
{\small\color{kursunkalem}(Solda küçük bir not: ``\okunamadi{iki değerin}'', iki akım için ok.)}
\begin{align*}
  I_1 &= I_S\,e^{V_1/V_T} \\
  I_2 &= I_S\,e^{V_2/V_T}
\end{align*}
\[ \frac{I_2}{I_1} = e^{(V_2 - V_1)/V_T} \]
\[ V_2 - V_1 = V_T\,\ln\frac{I_2}{I_1} \]
\[ V_2 - V_1 = (\ln 10\; V_T)\,\log\frac{I_2}{I_1} = \Big(\underbrace{2{,}3\,V_T}_{60\ \mathrm{mV}}\Big)\log\frac{I_2}{I_1} \]
{\small\color{kursunkalem}(Soluk, küçük not: ``\okunamadi{Wafer, IC process, IC package}''.)}

\medskip
\noindent \fbox{$i \simeq -I_S$} $\rightarrow$ Reverse Region

\begin{ornek}[4.4]
  \sekilnotu[19]{Devre: $V_{DD} = 5$ V pil (solda), $R = 1\,\mathrm{k}\Omega$ üstte, sağda aşağı bakan diyot; diyot üzerinde $V_D$ ($+$ üstte, $-$ altta), akım oku $I_D$.}
  The diode current is $1$ mA @ $0{,}7$ V
\end{ornek}
\begin{cozum}
  \noindent \fbox{\parbox{0.9\linewidth}{%
  ideal diyot modeli\\
  $V_D = 0\ \mathrm{V}$\\
  $\displaystyle I_D = \frac{5-0}{1} = 5\ \mathrm{mA}$}}

  \medskip
  \noindent \fbox{\parbox{0.9\linewidth}{%
  constant Voltage drop model (sabit gerilim düşümü modeli)\\
  $V_D = 0{,}7\ \mathrm{V}$\\
  $\displaystyle I_D = \frac{5-0{,}7}{1} = 4{,}3\ \mathrm{mA}$}}
\end{cozum}

\sayfa{20}
\section{Exponential model}
\[ I_D = \frac{V_{DD} - V_D}{R} \quad \textcircled{1} \qquad\qquad V_2 - V_1 = 2{,}3\,V_T\,\log\frac{I_2}{I_1} \quad \textcircled{2} \]
\[ I_D = \frac{5 - 0{,}7}{1} = 4{,}3\ \mathrm{mA} \]
\[ V_2 = V_1 + 0{,}06\,\log\frac{I_2}{I_1} \;\Longrightarrow\; V_2 = 0{,}738\ (\mathrm{V}) \]
\noindent $V_1 = 0{,}7$ V, \quad $I_1 = 1$ mA (ok ile $I_1$'e), \quad $I_2 = 4{,}3$ mA

\medskip
\noindent \underline{\textcolor{kirmizikalem}{Iteration 2nd}}
\begin{align*}
  I_D &= \frac{5 - 0{,}738}{1} = \kutula{4{,}262\ \mathrm{mA}} \\
  V_2 &= 0{,}738 + 0{,}06\left[\frac{4{,}262}{4{,}3}\right] \\
      &= \kutula{0{,}738\ \mathrm{V}}
\end{align*}
\footnote{Claude notu: muhtemelen $0{,}06\,\log\!\left[\frac{4{,}262}{4{,}3}\right]$ olmalı (sonuç $\approx 0{,}738$ V).}
\noindent {\small\color{kursunkalem}(Sağda, $V_2 = 0{,}738$ sonucundan kavisli bir çizgiyle bağlanmış not:)} ``Virgülden sonra 3 basamak aynı olduğundan iterasyonu durduruz''

\end{document}
